\documentclass[../DoAn.tex]{subfiles}
\begin{document}

\begin{center}
    \Large{\textbf{TÓM TẮT NỘI DUNG ĐỒ ÁN}}\\
\end{center}
\vspace{1cm}
Trong bối cảnh thị trường lao động ngày càng cạnh tranh, hồ sơ việc làm (CV) đóng vai trò then chốt trong quá trình ứng tuyển. Phần lớn các doanh nghiệp hiện nay sử dụng hệ thống theo dõi ứng viên (Applicant Tracking System -- ATS) để tự động sàng lọc CV trước khi đến tay nhà tuyển dụng. Tuy nhiên, nhiều ứng viên, đặc biệt là sinh viên mới ra trường, chưa nắm rõ cách tối ưu CV cho các hệ thống này. Các công cụ xây dựng CV trực tuyến hiện có chủ yếu cung cấp mẫu định dạng mà thiếu khả năng phân tích nội dung thông minh, đồng thời hầu hết chưa hỗ trợ tốt cho tiếng Việt.

Khóa luận này tiếp cận vấn đề trên bằng cách ứng dụng Mô hình Ngôn ngữ lớn (Large Language Model -- LLM), cụ thể là OpenAI GPT, kết hợp với nền tảng web hiện đại . Hệ thống cho phép người dùng tải lên CV dạng PDF hoặc DOCX, tự động trích xuất và phân tích nội dung bằng phương pháp lai (hybrid) kết hợp giữa kỹ thuật xử lý văn bản truyền thống và LLM, sau đó hỗ trợ chỉnh sửa trực quan theo từng phần thông qua trình soạn thảo có hướng dẫn.

Hệ thống được xây dựng trên nền tảng Next.js 15 với React 19, sử dụng TypeScript, tích hợp Supabase làm cơ sở dữ liệu và Vercel Blob làm dịch vụ lưu trữ tệp. Kết quả đạt được bao gồm: (i) pipeline phân tích CV hybrid giúp tiết kiệm token API đồng thời duy trì độ chính xác cao, (ii) trình chỉnh sửa CV tương tác hỗ trợ kéo thả và tự động lưu, (iii) khả năng xuất CV đa định dạng (PDF, DOCX, LaTeX, TXT), và (iv) hệ thống xác thực đa phương thức hỗ trợ song ngữ Việt--Anh.

\begin{flushright}
Sinh viên thực hiện\\
\begin{tabular}{@{}c@{}}
\textit{(Ký và ghi rõ họ tên)}\\[1.5cm]
Lê Sỹ Đan
\end{tabular}
\end{flushright}

\end{document}
